\documentclass[12pt]{article}
\usepackage{amsmath,amssymb,amsfonts,amsthm}
\usepackage[margin=1in]{geometry}

\begin{document}

\begin{center}
{\Large\bfseries Chapter 9, Problem 23}\\[6pt]
Byron \& Fuller, \textit{Mathematics of Classical and Quantum Physics}
\end{center}

\bigskip

\textbf{Problem.} Calculate the second Born approximation to the $s$-wave scattering amplitude, $F_0(p_1)$, defined by Eqs.~(9.112) and (9.113). Compare this result with the curves of Figs.~9.4, 9.5, and 9.6 and make any comments which seem appropriate.

\bigskip
\textbf{Solution.}

The $s$-wave ($l = 0$) scattering amplitude in the notation of the text (Eqs.~9.112, 9.113) is related to the partial-wave amplitude via the integral equation. The first Born approximation gives a real $F_0^{(1)}$.

The second Born approximation adds the term involving the iterated kernel:
\[
F_0^{(2)}(p_1) = F_0^{(1)}(p_1) + \Delta F_0(p_1),
\]
where $\Delta F_0$ is the second-order correction.

For the Yukawa potential $V(r) = g^2 e^{-\mu r}/r$ with dimensionless variables $p_1 = q_1/\mu$ (momentum in units of inverse range), the $s$-wave kernel involves $j_0(x) = \sin x / x$, so:

\textbf{First Born:}
\[
F_0^{(1)}(p_1) = -\frac{2m g^2}{\hbar^2 \mu}\int_0^\infty \frac{\sin^2(p_1 u)}{p_1^2 u}\,e^{-u}\,du,
\]
where $u = \mu r$.

This integral evaluates to:
\[
F_0^{(1)} = -\frac{mV_0}{\hbar^2\mu}\cdot\frac{1}{2p_1^2}\ln(1 + 4p_1^2),
\]
using the standard result $\int_0^\infty \frac{\sin^2(ax)}{x}\,e^{-x}\,dx = \frac{1}{4}\ln(1+4a^2)$.

\textbf{Second Born:} The correction involves a double integral with the Green's function $G_0(r,r') = -q_1\,r_<\,[n_0(q_1 r_>) + i j_0(q_1 r_>)]/r\,r'$ (for $l=0$, $j_0(x) = \sin x/x$, $n_0(x) = -\cos x/x$):
\[
\Delta F_0 = \left(\frac{2mg^2}{\hbar^2}\right)^2\!\!\int_0^\infty\!\!\int_0^\infty \frac{\sin(p_1 u)}{p_1 u}\,e^{-u}\left[\frac{e^{ip_1 u_>}\sin(p_1 u_<)}{p_1}\right]\frac{\sin(p_1 v)}{p_1 v}\,e^{-v}\,du\,dv.
\]

The imaginary part satisfies the optical theorem (Problem~22):
\[
\mathrm{Im}\,F_0^{(2)} = p_1|F_0^{(1)}|^2.
\]

The real part involves the principal-value integral with $n_0 = -\cos(q_1 r)/q_1 r$. This can be evaluated numerically for specific values of $p_1$.

\textbf{Comparison with figures:}

\begin{itemize}
\item \textbf{Fig.~9.4} ($mV_0/p^2\hbar^2 = 1.4$): The second Born approximation gives both real and imaginary parts of $F_0$. For small $p_1$, the Born series converges well and the second Born correction is small. For larger $p_1$, the correction becomes more significant but the series still converges because $\|K\| < 1$.

\item \textbf{Fig.~9.5} (convergence in the complex plane): The second Born approximation shifts the ``vector'' $F_0$ in the complex plane, with the imaginary part growing as $p_1$ increases. This is consistent with the spiral convergence pattern shown.

\item \textbf{Fig.~9.6} (weaker potential $mV_0/p^2\hbar^2 = 0.2$): For this weaker potential, the Born series converges much more rapidly, and the second Born approximation is already very close to the exact answer. The real and imaginary parts of $F_0^{(2)}$ nearly overlap with the exact curves.
\end{itemize}

The key observation is that the Born series converges better for weaker potentials and higher energies (larger $p_1$), and the optical theorem constraint $\mathrm{Im}\,F_0 = p_1|F_0|^2$ provides an important consistency check on the approximation. \hfill $\blacksquare$

\end{document}
